\chapter{Literature Review and Background}
\label{ch:background}

\section{Ising Formulations and Optimization}

Lucas (2014) proved that many NP-hard problems can be encoded as Ising Hamiltonians. An Ising model is a system of binary spins with interactions:
\[
H = \sum_i h_i s_i + \sum_{i<j} J_{ij} s_i s_j, \quad s_i \in \{-1, +1\}
\]

Finding the ground state (minimum energy) solves the optimization problem.

\section{Oscillator Ising Machines}

Wang and Roychowdhury (2019, 2021) developed OIM theory for coupled oscillator networks. Key papers:
\begin{itemize}
\item OIM dynamics: CMOS-compatible analog circuits
\item Multi-start convergence: 37--60\% success on dense graphs
\item Scalability: up to 100+ nodes with pruning
\end{itemize}

Recent advances use VO\textsubscript{2} oscillators, reaching nanosecond switching speeds.

\section{Multi-Robot Task Allocation}

Gerkey and Matarić (2004) established the MRTA taxonomy. Coalition formation (our focus) is NP-hard. Prior work:
\begin{itemize}
\item Exact solvers: exponential time, $<$20 agents
\item Greedy: O(N$^2$ M), near-optimal on trees
\item Distributed auction: asynchronous convergence
\end{itemize}

Gap: No prior work applies OIM to MRTA.

\section{Model Predictive Control}

MPC solves a constrained QP at each timestep. Standard approach:
\[
\min_u \|x - x_*\|^2 + \|u\|^2 \quad \text{subject to} \quad Ax + Bu = x_{next}
\]

Rawlings, Mayne \& Diehl (2020) provide the MPC textbook. For robots, standard solvers (OSQP, Mosek) require milliseconds per solve on embedded hardware.

Gap: No prior mapping of PIPG to SNNs.

\section{Neuromorphic Computing Platforms}

\begin{itemize}
\item Intel Loihi: Spiking neural cores, event-driven
\item IBM TrueNorth: 5.4B transistors, 1 mW
\item Brainscales-2: Analog neurons, configurable plasticity
\end{itemize}

Mangalore et al. (2024) first demonstrated MPC on Loihi. Our work extends this framework.

